\unitopen{38}
\begin{candos}
\item rozpoznawać i rozwijać imiesłowowe określenia przed rzeczownikiem (\de{die steigenden Kosten, der entwickelte Dienst});
\item tworzyć rozszerzone określenia (\de{die im Test verwendeten Daten});
\item zamieniać je na zdania względne i odwrotnie;
\item czytać tekst specjalistyczny z długimi grupami przed rzeczownikiem.
\end{candos}

\section{Słownictwo}
\vocabtable{38}

\section{Imiesłowy jako przymiotniki}
\textbf{Partizip I} (\de{-d}): czynność trwa, strona czynna: \de{laufen} $\rightarrow$ \de{das laufende Projekt}, \de{steigen} $\rightarrow$
\de{die steigenden Kosten}. \textbf{Partizip II}: czynność zakończona albo strona bierna: \de{entwickeln} $\rightarrow$ \de{der entwickelte
Dienst} (opracowany), \de{speichern} $\rightarrow$ \de{die gespeicherten Daten}. Dostają końcówki przymiotnika (jednostki 21 i 27).

\begin{center}\small
\begin{tabular}{@{}p{0.45\linewidth}p{0.47\linewidth}@{}}
\toprule
\textsf{\textbf{Zdanie względne}} & \textsf{\textbf{Określenie imiesłowowe}}\\
\midrule
\de{die Kosten, die steigen} & \de{die steigenden Kosten}\\
\de{der Dienst, der entwickelt wurde} & \de{der entwickelte Dienst}\\
\de{die Daten, die gespeichert sind} & \de{die gespeicherten Daten}\\
\de{die Daten, die im Test verwendet wurden} & \de{die im Test verwendeten Daten}\\
\de{der Dienst, den wir entwickelt haben} & \de{der von uns entwickelte Dienst}\\
\bottomrule
\end{tabular}
\end{center}

\begin{zasada}[Rozszerzone określenie]
Grupa wstawiana między rodzajnik a rzeczownik: \de{der} + [okoliczniki, dopełnienia] + \emph{imiesłów z końcówką} + rzeczownik.
Czytaj \emph{od rodzajnika do rzeczownika}, znajdź imiesłów na końcu grupy i potraktuj wszystko przed nim jako jego dopełnienia:
\de{die [in der letzten Woche vom Team getesteten] Funktionen}.
\end{zasada}

\begin{falle}
Polski przekłada takie konstrukcje zdaniem względnym („dane użyte w teście”), a niemiecki umieszcza je \emph{przed} rzeczownikiem. Błąd
typowy: pisanie \de{die Daten verwendet im Test} (szyk polski). Drugi: pomylone Partizip I i II (\de{laufend} to trwa, \de{gelaufen} to zakończone).
Trzeci: brak końcówki przymiotnika po imiesłowie przed rzeczownikiem.
\end{falle}

\begin{szwajcaria}
Szwajcarskie raporty roczne i dokumenty prawne są pełne rozszerzonych określeń. W pracy zawodowej ich używaj oszczędnie, ale ucz się
je rozumieć: skraca czas czytania dokumentacji technicznej, specyfikacji i umów.
\end{szwajcaria}

\begin{praca}
\de{Die im letzten Sprint entwickelten Funktionen sind getestet.} \de{Die zu erwartenden Kosten steigen.} \de{Der von
allen Teams genutzte Dienst ist kritisch.} Przepisz opis swojego PR tak, żeby użyć dwóch takich określeń.
\end{praca}

\section{Dialogi}
\textbf{Dialog 1 --- opis wyników}
\begin{dialog}
\turn{A}{Was zeigen die im Test gemessenen Werte?}{Co pokazują wartości zmierzone w teście?}
\turn{B}{Die steigenden Kosten hängen mit dem wachsenden Datenvolumen zusammen.}{Rosnące koszty wiążą się z rosnącą ilością danych.}
\turn{A}{Und der von uns entwickelte Cache?}{A opracowany przez nas cache?}
\turn{B}{Er verkürzt die laufenden Abfragen um etwa ein Drittel.}{Skraca trwające zapytania o około jedną trzecią.}
\end{dialog}

\section{Ćwiczenia}
\begin{cwiczenia}
\exitem{Partizip I jako przymiotnik: \de{steigen} (die \dots{} Kosten).}{\de{die steigenden Kosten}}
\exitem{Partizip II jako przymiotnik: \de{speichern} (die \dots{} Daten).}{\de{die gespeicherten Daten}}
\exitem{Zamień na określenie: \de{die Daten, die im Test verwendet wurden}.}{\de{die im Test verwendeten Daten}}
\exitem{Zamień na określenie: \de{der Dienst, den wir entwickelt haben}.}{\de{der von uns entwickelte Dienst}}
\exitem{Zamień na zdanie względne: \de{die laufende Anfrage}.}{\de{die Anfrage, die läuft}}
\exitem{Zamień na zdanie względne: \de{der geplante Termin}.}{\de{der Termin, der geplant ist}}
\exitem{Wskaż Partizip I: \de{die laufenden Kosten} czy \de{die gelaufenen Kosten}?}{\de{laufenden}}
\exitem{Popraw: \de{die Daten verwendet im Test}.}{\de{die im Test verwendeten Daten}}
\exitem{Przetłumacz: „Funkcje opracowane w ostatnim sprincie są przetestowane.”}{\de{Die im letzten Sprint entwickelten Funktionen sind getestet.}}
\exitem{Przetłumacz: „Rosnące ryzyko wymaga decyzji.”}{\de{Das steigende Risiko erfordert eine Entscheidung.}}
\end{cwiczenia}

\weekend{Wybierz techniczny akapit (dokumentacja, artykuł) i znajdź w nim pięć określeń imiesłowowych. Rozbij każde na zdanie względne.
Napisz własny akapit o projekcie z trzema takimi określeniami. Nagraj streszczenie.}
